\section{Problemas y objetivos}

El problema general es la ausencia de una cadena local y reproducible que permita pasar de subtítulos públicos a alertas temporales comparables y, después, a una decisión humana informada. Cuatro problemas específicos vinculan esa situación con las acciones realizadas y con la evidencia que permite comprobar el cambio alcanzado.

Primero, el proyecto no cuenta con un corpus local versionado que conserve el texto, el intervalo temporal, la fuente, las categorías y la procedencia de cada etiqueta. Esta ausencia impide reconstruir qué observa cada modelo y comprobar que un mismo video no aparezca en entrenamiento y prueba. La definición del daño, el muestreo y la anotación condicionan el comportamiento posterior de los clasificadores de lenguaje abusivo \cite{vidgen2020directions}; además, los \emph{data statements} y las \emph{data cards} recomiendan documentar origen, recolección, anotación, usos y límites \cite{bender2018datastatements,pushkarna2022datacards}. Por ello se construye, depura y documenta un corpus de subtítulos públicos con particiones agrupadas por video. El corpus integrado, su manifiesto y la taxonomía verifican el artefacto de datos dentro del alcance del estudio.

Segundo, los resultados disponibles proceden de etapas y contratos distintos, por lo que no bastan para decidir qué familia o estructura separa mejor las categorías de moderación poco frecuentes. En conjuntos desbalanceados, las métricas agregadas pueden ocultar el desempeño de la clase positiva, y las curvas precisión--recobrado son más informativas para ese análisis que las curvas ROC \cite{sokolova2009metrics,saito2015pr}. El objetivo consiste en entrenar y comparar modelos clásicos, MiniLM/E5 y el LLM Qwen3--0.6B en variantes planas, en cascada y jerárquicas, con una partición y un contrato comunes. La comparación común reduce la incertidumbre entre enfoques dentro de ese protocolo; confirmar la generalización requiere nuevas semillas y una cohorte temporal \cite{dror2018significance,cawley2010selection}.

Tercero, una etiqueta obligatoria oculta los casos dudosos y no muestra cuántos fragmentos debe revisar una persona ni cuántos daños quedan en la ruta automática. La clasificación selectiva permite abstenerse ante casos inciertos, y el aprendizaje para diferir extiende esa decisión al envío a una persona experta \cite{chow1970reject,geifman2017selective,mozannar2020defer}. Por eso se establece como objetivo calibrar los scores solo con validación y comparar candidatos al mismo recobrado de enrutamiento del 95\%, con precisión, tasa de revisión y falsos negativos. La política congelada, sus umbrales y la evaluación en test hacen verificable esa ruta de decisión: el sistema cuantifica la incertidumbre, los falsos negativos y la carga humana.

Cuarto, los modelos, las métricas y los archivos permanecen separados y no forman un artefacto que un supervisor pueda utilizar ni una memoria preparada para correcciones posteriores. DSR vincula la construcción de un artefacto con su demostración y evaluación iterativa \cite{hevner2004dsr,peffers2007dsrm}; en moderación, además, la predicción automática no sustituye las reglas, el trabajo humano ni la responsabilidad sobre la decisión \cite{gorwa2020moderation}. El objetivo consiste en empaquetar los modelos seleccionados, la evidencia temporal, el consenso, la revisión humana y el registro estadístico en un mismo flujo. El frontend 05, el registro de modelos, la base de datos y la ontología muestran que el prototipo integra los resultados en un flujo semiautomático trazable. Así se resuelve el aislamiento técnico en el nivel de demostrador; su eficacia bajo tráfico real todavía requiere un piloto supervisado.

La Tabla~\ref{tab:problemas} resume esta relación entre los cuatro problemas, las acciones realizadas y la evidencia que permite comprobarlas.

\begin{table*}[t]
\centering
\caption{Síntesis de problemas, objetivos y evidencia de verificación}
\label{tab:problemas}
\small
\begin{tabularx}{\textwidth}{>{\bfseries}p{12mm}X X X}
\toprule
 & \textbf{Problema específico} & \textbf{Objetivo específico} & \textbf{Evidencia de verificación} \\
\midrule
P1/O1 & No existe un corpus local versionado y separado por video. & Construir y documentar el corpus integrado. & 117\,244 chunks, 3\,230 videos, manifiesto, taxonomía y particiones agrupadas. \\
P2/O2 & Falta una comparación común entre familias y estructuras. & Comparar clásicos, MiniLM/E5 y Qwen en diseños planos y estructurados. & 14 configuraciones evaluadas sobre el mismo test activo de cuatro categorías. \\
P3/O3 & Las decisiones forzadas ocultan incertidumbre, carga y omisiones. & Calibrar en validación y comparar al 95\% de recobrado de enrutamiento. & Umbrales congelados; precisión, revisión y falsos negativos reportados. \\
P4/O4 & Los resultados no forman un flujo utilizable y reentrenable. & Integrar modelos, evidencia, consenso, revisión y registro. & Frontend 05, registro de modelos, base de datos y ontología. \\
\bottomrule
\end{tabularx}
\end{table*}
